\subsection{解析半群}\label{sec:analytic-semigroups}

\subsubsection{解析半群及其生成元}\label{subsec:analytic-semigroups-generators}

现在讨论一类半群, 它在许多方面可得到比强连续半群强得多的结果. 特别地, 将得到:\begin{enumerate}\item 初值问题的解具有更好的正则性.\item 关于无穷小生成元扰动的更好结果.\item 半群的型与无穷小生成元的谱之间的关系.\end{enumerate}先给出解析半群的定义.

\begin{Definition}
\label{def:analytic-semigroup}
若强连续半群 $\exp(At)$ 满足以下条件, 则称其为\textbf{解析半群}:\begin{enumerate}\item 对某个 $\theta\in(0,\pi/2)$, $\exp(At)\in\mathcal{L}(X)$ 可延拓到 $t\in\Delta_\theta=\{0\}\cup\{t\in\mathbb{C}\mid|\arg t|<\theta\}$, 且 Definition~\ref{def:strongly-continuous-semigroup} 的条件对 $t\in\Delta_\theta$ 成立.\item 对 $t\in\Delta_\theta\setminus\{0\}$, $\exp(At)$ 关于 $t$ 解析 (按一致算子拓扑的意义).\end{enumerate}
\end{Definition}

现在最重要的问题是如何刻画解析半群的无穷小生成元. 下面的定理处理这一问题.

\begin{Theorem}
\label{thm:analytic-generator-sectorial-resolvent}
$X$ 中闭且稠密定义的算子 $A$ 为解析半群的生成元, 当且仅当存在 $\omega\in\mathbb{R}$, 使半平面 $\operatorname{Re}\lambda>\omega$ 包含于 $A$ 的预解集, 并且存在常数 $C$, 使
\begin{equation}\label{eq:analytic-semigroup-half-plane-resolvent-estimate}
\|R_\lambda(A)\|\leq C/|\lambda-\omega|
\end{equation}
对 $\operatorname{Re}\lambda>\omega$ 成立. 若如此, 则实际上预解集对某个 $\delta>0$ 包含扇形域 $|\arg(\lambda-\omega)|<\frac{\pi}{2}+\delta$, 且该扇形域中有类似的预解算子估计. 此外, 半群可表示为
\begin{equation}\label{eq:analytic-semigroup-contour-integral-representation}
e^{At}=\dfrac{1}{2\pi i}\int_\Gamma e^{\lambda t}(\lambda I-A)^{-1}\,\mathrm{d}\lambda,
\end{equation}
其中 $\Gamma$ 为从 $e^{-i\phi}\infty$ 到 $e^{i\phi}\infty$ 的任意曲线, 且 $\Gamma$ 完全位于集合 $\{|\arg(\lambda-\omega)|\leq\phi\}$ 中. 这里 $\phi$ 为满足 $\frac{\pi}{2}<\phi<\frac{\pi}{2}+\delta$ 的任意角.
\end{Theorem}

\begingroup
\renewcommand{\endProof}{\par}
\begin{Proof}
条件的必要性不难看出. 假设 $\exp(At)$ 为定义在 $\Delta_\theta$ 中的解析半群. 容易证明, 对任意 $\psi<\theta$, 都存在常数 $M$ 和 $\omega$, 使对 $t\in\Delta_\psi$ 有 $\|\exp(At)\|\leq M\exp(\omega|t|)$; 证明沿用 Theorem~\ref{thm:semigroup-logarithmic-growth-type} 的思路. 由 Hille-Yosida 定理的证明可知
\begin{equation}\label{eq:analytic-generator-necessity-laplace-resolvent}
-R_\lambda(A)=\int_0^\infty e^{-\lambda t}T(t)\,\mathrm{d}t
\end{equation}
对 $\operatorname{Re}\lambda>\omega$ 成立. 不失一般性, 设 $\operatorname{Im}\lambda>0$. 然后将积分路径移至射线 $\arg t=-\delta$, $\delta=\max(\psi,\arg\lambda)$, 得到
\begin{equation}\label{eq:analytic-generator-rotated-laplace-integral}
-R_\lambda(A)=\int_0^\infty\exp(-\lambda e^{-i\delta}t)T(e^{-i\delta}t)e^{-i\delta}\,\mathrm{d}t,
\end{equation}
从而
\begin{equation}\label{eq:analytic-generator-rotated-resolvent-bound}
\|R_\lambda(A)\|\leq\dfrac{M}{\operatorname{Re}(\lambda\exp(-i\delta))-\omega}\leq\dfrac{C}{|\lambda-\omega|}.
\end{equation}

\end{Proof}
\endgroup

现在假设 \eqref{eq:analytic-semigroup-half-plane-resolvent-estimate} 对 $\operatorname{Re}\lambda>\omega$ 成立. 若 $\lambda-\omega=Re^{i\theta}$, $\theta>\pi/2$, 且 $\phi$ 为任意小于 $\pi/2$ 的角, 则可写成
\begin{equation}\label{eq:analytic-generator-sector-resolvent-extension-splitting}
A-\lambda I=A-(\omega+Re^{i\phi})I-R(e^{i\theta}-e^{i\phi})I,
\end{equation}
只要 $|e^{i\theta}-e^{i\phi}|\leq C$, 其逆就可表示为几何级数. 这说明预解集延伸到超出半平面 $\operatorname{Re}\lambda>\omega$ 的某个扇形域, 而对几何级数逐项估计可得到预解算子估计. 特别地, 这使 \eqref{eq:analytic-semigroup-contour-integral-representation} 有良好定义.

接下来证明 \eqref{eq:analytic-semigroup-contour-integral-representation} 确实表示由 $A$ 生成的解析半群. 显然, \eqref{eq:analytic-semigroup-contour-integral-representation} 对 $|\arg t|<\phi-\frac{\pi}{2}$ 有定义, 且关于 $t$ 解析. 再验证半群性质. 用 $T(t)$ 表示 \eqref{eq:analytic-semigroup-contour-integral-representation} 的右端, 取 $\Gamma^{\prime}$ 为具有与 $\Gamma$ 相同性质但完全位于 $\Gamma$ 右侧的路径. 则有\[\begin{aligned}
T(s)T(t)&=\left(\dfrac{1}{2\pi i}\right)^2\int_{\Gamma^{\prime}}\int_\Gamma e^{\lambda^{\prime}s+\lambda t}(\lambda^{\prime}I-A)^{-1}(\lambda I-A)^{-1}\,\mathrm{d}\lambda\,\mathrm{d}\lambda^{\prime}\\
&=\left(\dfrac{1}{2\pi i}\right)^2\Bigg[\int_{\Gamma^{\prime}}e^{\lambda^{\prime}s}(\lambda^{\prime}I-A)^{-1}\int_\Gamma e^{\lambda t}(\lambda-\lambda^{\prime})^{-1}\,\mathrm{d}\lambda\,\mathrm{d}\lambda^{\prime}\\
&\qquad-\int_\Gamma e^{\lambda t}(\lambda I-A)^{-1}\int_{\Gamma^{\prime}}e^{\lambda^{\prime}s}(\lambda-\lambda^{\prime})^{-1}\,\mathrm{d}\lambda^{\prime}\,\mathrm{d}\lambda\Bigg]\\
&=\dfrac{1}{2\pi i}\int_\Gamma e^{\lambda(t+s)}(\lambda I-A)^{-1}\,\mathrm{d}\lambda=T(t+s).
\end{aligned}\]这里使用了恒等式 $(\lambda I-A)^{-1}(\lambda^{\prime}I-A)^{-1}=(\lambda-\lambda^{\prime})^{-1}[(\lambda^{\prime}I-A)^{-1}-(\lambda I-A)^{-1}]$ 以及事实\[\begin{aligned}
\int_\Gamma e^{\lambda t}(\lambda-\lambda^{\prime})^{-1}\,\mathrm{d}\lambda&=0,\\
\int_{\Gamma^{\prime}}e^{\lambda^{\prime}s}(\lambda-\lambda^{\prime})^{-1}\,\mathrm{d}\lambda^{\prime}&=-2\pi i e^{\lambda s}.
\end{aligned}\]

必要时将积分路径右移, 可假设 \eqref{eq:analytic-semigroup-contour-integral-representation} 中的 $\Gamma$ 完全位于集合 $\{|\arg\lambda|<\phi\}$ 中. 对足够小的 $t$, 可作代换 $\lambda^{\prime}=\lambda|t|$, 再将积分路径变形回 $\Gamma$. 由此得到
\begin{equation}\label{eq:analytic-semigroup-contour-rescaled-at-small-time}
T(t)=\dfrac{1}{2\pi i|t|}\int_\Gamma e^{\lambda^{\prime}t/|t|}\left(\dfrac{\lambda^{\prime}}{|t|}I-A\right)^{-1}\,\mathrm{d}\lambda^{\prime}.
\end{equation}
从而得到
\begin{equation}\label{eq:analytic-semigroup-small-time-contour-norm-bound}
\|T(t)\|\leq\dfrac{C}{2\pi}\int_\Gamma|e^{\lambda^{\prime}t/|t|}||\lambda^{\prime}|^{-1}\,\mathrm{d}\lambda^{\prime}.
\end{equation}
显然, 当 $t\to0$ 时, 该表达式保持有界. 此外, 对 $u\in\mathcal{D}(A)$, 若选择 $\Gamma$ 位于原点右侧, 则得到\[\begin{aligned}
T(t)u-u&=\dfrac{1}{2\pi i}\int_\Gamma e^{\lambda t}((\lambda I-A)^{-1}-\lambda^{-1})u\,\mathrm{d}\lambda\\
&=-\dfrac{1}{2\pi i}\int_\Gamma e^{\lambda t}\lambda^{-1}(\lambda I-A)^{-1}\,\mathrm{d}\lambda\ Au.
\end{aligned}\]当 $t\to0$ 时, 该积分收敛到
\begin{equation}\label{eq:analytic-semigroup-origin-contour-zero-integral}
-\dfrac{1}{2\pi i}\int_\Gamma\lambda^{-1}(\lambda I-A)^{-1}\,\mathrm{d}\lambda=0,
\end{equation}
这可通过在右侧用圆弧闭合积分路径看出. 因而 $T(t)$ 在 $t\to0$ 时强收敛到恒等算子.

\begingroup
\renewenvironment{Proof}{\par}{\par\nobreak\hfill$\blacksquare$\par\vspace{0.1cm}}
\begin{Proof}
对 $t\neq0$, 由 \eqref{eq:analytic-semigroup-contour-integral-representation} 得到:\[\begin{aligned}
\dfrac{\mathrm{d}}{\mathrm{d}t}T(t)&=\dfrac{1}{2\pi i}\int_\Gamma\lambda e^{\lambda t}(\lambda I-A)^{-1}\,\mathrm{d}\lambda\\
&=\dfrac{1}{2\pi i}\int_\Gamma e^{\lambda t}A(\lambda I-A)^{-1}\,\mathrm{d}\lambda+\dfrac{1}{2\pi i}\int_\Gamma e^{\lambda t}I\,\mathrm{d}\lambda\\
&=AT(t).
\end{aligned}\]这里利用 $A$ 的闭性来说明可将其提出积分号. 因而对 $u\in\mathcal{D}(A)$ 有\[\begin{aligned}
\dfrac{T(h)-I}{h}u&=\lim_{\epsilon\to0}\dfrac{T(h)-T(\epsilon)}{h}u\\
&=\lim_{\epsilon\to0}\dfrac{1}{h}\int_\epsilon^h AT(t)u\,\mathrm{d}t=\dfrac{1}{h}\int_0^h T(t)Au
\end{aligned}\]当 $h\to0$ 时, 极限为 $Au$; 因而该半群的无穷小生成元是 $A$ 的延拓. 由于 $A$ 在半平面 $\operatorname{Re}\lambda>\omega$ 中有预解算子, 这一延拓不可能是真延拓.
\end{Proof}
\endgroup

在上述证明中, 已证明 $\exp(At)$ 对 $t>0$ 可微, 且 $\frac{\mathrm{d}}{\mathrm{d}t}\exp(At)=A\exp(At)$. 采用本质上相同的论证, 可证明 $\exp(At)$ 实际上无穷次可微, 且 $\frac{\mathrm{d}^n}{\mathrm{d}t^n}\exp(At)=A^n\exp(At)$. 特别地, 这意味着对每个 $n$, $\exp(At)$ 的值域都包含于 $\mathcal{D}(A^n)$. 在 PDE 应用中, 这通常转化为光滑化性质. 即使初始条件具有奇性, 解对正的 $t$ 也是光滑的 (为 $C^\infty$ 类). 为以后使用, 记录下面对 $A\exp(At)$ 的范数的界.

\begin{Lemma}
\label{lem:analytic-semigroup-time-derivative-bound}
设 $A$ 为解析半群的无穷小生成元. 则存在常数 $C$ 和 $\omega$, 使对每个 $t>0$ 都有
\begin{equation}\label{eq:analytic-semigroup-generator-time-derivative-bound}
\|A\exp(At)\|\leq C\exp(\omega t)/t.
\end{equation}

\end{Lemma}

证明基于与导出 \eqref{eq:analytic-semigroup-small-time-contour-norm-bound} 类似的计算; 将细节留作习题 (Problem~\ref{ex:analytic-semigroup-derivative-bound-proof}).

在 ODE 课程中学到, 若 $A$ 的所有特征值都具有负实部, 则方程 $\dot u=Au$ 的零解稳定. 实际上, 这是评估稳定性时应用最广泛的实用判据. 自然希望将同类判据用于 PDE 所描述的物理系统. 这提出以下问题: 若 $A$ 为某半群的生成元, 且 $A$ 的谱完全位于左半平面并与虚轴保持正距离, 是否可推出 $\|\exp(At)\|$ 在 $t\to\infty$ 时有界? 或等价地, $A$ 的谱的实部的上界与半群的型之间是否存在联系? 遗憾的是, 在一般 $C_0$ 半群的情形下, 答案是否定的. 虽然反例并非物理应用中会 ``自然\textquotesingle\textquotesingle\ 出现的那一类, 但尚无已知的实用条件可得到肯定结果. 目前的认识状况显然并不令人满意. 然而, 对解析半群, 这一问题很容易. 有以下结果.

\begin{Theorem}
\label{thm:analytic-semigroup-spectrum-controls-growth}
设 $A$ 为解析半群的无穷小生成元, 并假设 $A$ 的谱完全位于直线 $\operatorname{Re}\lambda=\omega$ 的左侧. 则存在常数 $M$, 使 $\|\exp(At)\|\leq M\exp(\omega t)$.
\end{Theorem}

证明的思路是移动 \eqref{eq:analytic-semigroup-contour-integral-representation} 中的积分路径, 使其完全位于 $\operatorname{Re}\lambda=\omega$ 的左侧. 将论证细节留作习题 (Problem~\ref{ex:analytic-semigroup-spectral-growth-proof}).

\subsubsection{分数次幂}\label{subsec:analytic-generator-fractional-powers}

本小节始终假设 $A$ 为解析半群的无穷小生成元, 且 $A$ 的谱完全位于 (开) 左半平面中. 在这种情形下, 将定义 $-A$ 的分数次幂. 若 $A$ 的谱位于半平面 $\operatorname{Re}\lambda<\omega$ 中, 当然也可用同样的考虑定义 $\omega I-A$ 的分数次幂. 若谱在开左半平面中, 可选取 $\delta>0$, 使 $A$ 的谱位于半平面 $\operatorname{Re}\lambda<-\delta$ 中; 由前一小节, 有如下形式的界 $\|\exp(At)\|\leq M\exp(-\delta t)$, $\|A\exp(At)\|\leq M_1\exp(-\delta t)/t$. 此外, 注意到 $A^n\exp(At)=(A\exp(A\frac{t}{n}))^n$, 得到 $\|A^n\exp(At)\|\leq M_n\exp(-\delta t)/t^n$.

定义
\begin{equation}\label{eq:analytic-generator-negative-fractional-power-contour}
(-A)^{-\alpha}=-\dfrac{1}{2\pi i}\int_\Gamma\lambda^{-\alpha}(\lambda I+A)^{-1}\,\mathrm{d}\lambda,
\end{equation}
其中 $\Gamma$ 为从 $\exp(-i\theta)\infty$ 到 $\exp(i\theta)\infty$ 的曲线, 使 $-A$ 的谱位于其右侧, 原点位于 $\Gamma$ 左侧. 这里 $\pi/2-\delta<\theta<\pi$, $\delta$ 如 Theorem~\ref{thm:analytic-generator-sectorial-resolvent} 中所给, 而 $\lambda^{-\alpha}$ 表示在正实轴上取正值的函数分支. 由 \eqref{eq:analytic-semigroup-half-plane-resolvent-estimate} 可知, 对任意 $\alpha>0$, \eqref{eq:analytic-generator-negative-fractional-power-contour} 中的积分都绝对收敛. 若 $\alpha$ 为整数, 可将 \eqref{eq:analytic-generator-negative-fractional-power-contour} 中的积分路径变形为绕原点的圆, 并利用留数计算该积分; 确实得到 $(-A)^{-n}=(-A^{-1})^n$; 见 Problem~\ref{ex:fractional-power-integer-inverse-identity}. 若 $0<\alpha<1$, 可将 $\Gamma$ 变形为负实轴的上侧和下侧. 这就得到表达式
\begin{equation}\label{eq:analytic-generator-negative-fractional-power-real-resolvent}
(-A)^{-\alpha}=\dfrac{\sin(\pi\alpha)}{\pi}\int_0^\infty\lambda^{-\alpha}(\lambda I-A)^{-1}\,\mathrm{d}\lambda,\ 0<\alpha<1.
\end{equation}


回顾 Hille-Yosida 定理的证明, 有
\begin{equation}\label{eq:fractional-power-resolvent-laplace-formula}
(\lambda I-A)^{-1}=\int_0^\infty e^{-\lambda t}e^{At}\,\mathrm{d}t.
\end{equation}
将这一表达式代入 \eqref{eq:analytic-generator-negative-fractional-power-real-resolvent}, 再交换积分次序, 得到\[\begin{aligned}
(-A)^{-\alpha}&=\dfrac{\sin(\pi\alpha)}{\pi}\int_0^\infty\exp(At)\int_0^\infty\lambda^{-\alpha}e^{-\lambda t}\,\mathrm{d}\lambda\,\mathrm{d}t\\
&=\dfrac{\sin(\pi\alpha)}{\pi}\left(\int_0^\infty u^{-\alpha}e^{-u}\,\mathrm{d}u\right)\left(\int_0^\infty t^{\alpha-1}e^{At}\,\mathrm{d}t\right).
\end{aligned}\]有
\begin{equation}\label{eq:fractional-power-gamma-reflection-integral}
\int_0^\infty u^{-\alpha}e^{-u}\,\mathrm{d}u=\dfrac{\pi}{\sin(\pi\alpha)\Gamma(\alpha)},
\end{equation}
因此
\begin{equation}\label{eq:analytic-generator-fractional-power-semigroup-integral}
(-A)^{-\alpha}=\dfrac{1}{\Gamma(\alpha)}\int_0^\infty t^{\alpha-1}e^{At}\,\mathrm{d}t.
\end{equation}
这里给出的论证仅在 $0<\alpha<1$ 时适用. 然而, \eqref{eq:analytic-generator-negative-fractional-power-contour} 和 \eqref{eq:analytic-generator-fractional-power-semigroup-integral} 对任意 $\alpha>0$ 都有定义, 且关于 $\alpha$ 解析, 因而由解析延拓的唯一性, 它们对每个 $\alpha>0$ 都相同. 利用 \eqref{eq:analytic-generator-fractional-power-semigroup-integral} 以及界 $\|\exp(At)\|\leq M\exp(-\delta t)$, 容易建立下面的结果.

\begin{Lemma}
\label{lem:negative-fractional-power-uniform-bound}
存在常数 $C$, 使对 $0<\alpha\leq1$ 有 $\|(-A)^{-\alpha}\|\leq C$.
\end{Lemma}

\begin{Theorem}
\label{thm:negative-fractional-powers-semigroup}
有以下结论.\begin{enumerate}\item $(-A)^{-\alpha}(-A)^{-\beta}=(-A)^{-(\alpha+\beta)}$.\item $\lim_{\alpha\to0}(-A)^{-\alpha}=I$ 在强算子拓扑下成立.\end{enumerate}
\end{Theorem}

采用显然的约定 $(-A)^0=I$, 这一定理断言算子 $(-A)^{-\alpha}$, $\alpha\geq0$, 构成 $C_0$ 半群. 自然可将无穷小生成元称为 $-\log(-A)$, 但这里不再深入讨论.

\begin{Proof}
有\[\begin{aligned}
&(-A)^{-\alpha}(-A)^{-\beta}\\
&\quad=\dfrac{1}{\Gamma(\alpha)\Gamma(\beta)}\int_0^\infty\int_0^\infty t^{\alpha-1}s^{\beta-1}e^{At}e^{As}\,\mathrm{d}t\,\mathrm{d}s\\
&\quad=\dfrac{1}{\Gamma(\alpha)\Gamma(\beta)}\int_0^\infty t^{\alpha-1}\int_t^\infty(u-t)^{\beta-1}e^{Au}\,\mathrm{d}u\,\mathrm{d}t\\
&\quad=\dfrac{1}{\Gamma(\alpha)\Gamma(\beta)}\int_0^\infty\int_0^u t^{\alpha-1}(u-t)^{\beta-1}\,\mathrm{d}t\ e^{Au}\,\mathrm{d}u\\
&\quad=\dfrac{1}{\Gamma(\alpha)\Gamma(\beta)}\int_0^1 v^{\alpha-1}(1-v)^{\beta-1}\,\mathrm{d}v\int_0^\infty u^{\alpha+\beta-1}e^{Au}\,\mathrm{d}u\\
&\quad=\dfrac{1}{\Gamma(\alpha+\beta)}\int_0^\infty u^{\alpha+\beta-1}e^{Au}\,\mathrm{d}u\\
&\quad=(-A)^{-\alpha-\beta}.
\end{aligned}\]由于已经知道 $\|(-A)^{-\alpha}\|$ 在 $\alpha\to0$ 时有界, 只需对 $X$ 的某个稠密子集中的 $u$ 证明 $(-A)^{-\alpha}u\to u$. 选取 $u\in\mathcal{D}(A)$, 则对某个 $y\in X$ 有 $u=-A^{-1}y$. 于是有
\begin{equation}\label{eq:fractional-power-semigroup-strong-continuity-reduction}
(-A)^{-\alpha}u-u=(-A)^{-1-\alpha}y-(-A)^{-1}y,
\end{equation}
由 \eqref{eq:analytic-generator-negative-fractional-power-contour} 或 \eqref{eq:analytic-generator-fractional-power-semigroup-integral} 显然可知, 当 $\alpha>0$ 时, $(-A)^{-\alpha}$ 在一致算子拓扑下实际上连续 (甚至解析).
\end{Proof}

由于 $(-A)^{-n}$ 对 $n\in\mathbb{N}$ 为单射, 且对 $n>\alpha$ 有 $(-A)^{-n}=(-A)^{-n+\alpha}(-A)^{-\alpha}$, 可知 $(-A)^{-\alpha}$ 为单射. 因而它有逆, 自然记为 $(-A)^\alpha$. 显然 $(-A)^\alpha$ 为闭算子, 其定义域为 $\mathcal{D}((-A)^\alpha)=\mathcal{R}((-A)^{-\alpha})$; 由于对 $n>\alpha$ 有 $\mathcal{R}((-A)^{-n})=\mathcal{D}(A^n)\subset\mathcal{R}((-A)^{-\alpha})$, 可知 $(-A)^\alpha$ 的定义域稠密. 此外, 容易验证, 对任意 $\alpha,\beta\in\mathbb{R}$ 和任意 $u\in\mathcal{D}((-A)^\gamma)$, 都有 $(-A)^{\alpha+\beta}u=(-A)^\alpha(-A)^\beta u$, 其中 $\gamma=\max(\alpha,\beta,\alpha+\beta)$.\par 以联系 $(-A)^\alpha$ 与半群的一个结果结束本小节.

\begin{Lemma}
\label{lem:fractional-power-semigroup-smoothing-holder}
设 $\alpha>0$. 对每个 $u\in\mathcal{D}((-A)^\alpha)$, 有\[\exp(At)(-A)^\alpha u=(-A)^\alpha\exp(At)u.\]此外, 算子 $(-A)^\alpha\exp(At)$ 有界, 且具有如下形式的界
\begin{equation}\label{eq:analytic-semigroup-fractional-power-smoothing-bound}
\|(-A)^\alpha\exp(At)\|\leq M_\alpha t^{-\alpha}e^{-\delta t}.
\end{equation}
若 $0<\alpha\leq1$ 且 $u\in\mathcal{D}((-A)^\alpha)$, 则有如下形式的界
\begin{equation}\label{eq:analytic-semigroup-fractional-domain-holder-orbit-bound}
\|\exp(At)u-u\|\leq C_\alpha t^\alpha\|(-A)^\alpha u\|.
\end{equation}

\end{Lemma}

\begin{Proof}
若 $u\in\mathcal{D}((-A)^\alpha)$, 则对某个 $v$ 有 $u=(-A)^{-\alpha}v$, 从而得到\[\begin{aligned}
\exp(At)u&=\exp(At)(-A)^{-\alpha}v=\dfrac{1}{\Gamma(\alpha)}\int_0^\infty s^{\alpha-1}\exp(As)\exp(At)v\,\mathrm{d}s\\
&=(-A)^{-\alpha}\exp(At)v=(-A)^{-\alpha}\exp(At)(-A)^\alpha u.
\end{aligned}\]对两端作用 $(-A)^\alpha$, 即得引理的第一个断言.\par 设 $n-1<\alpha<n$, 则\[\begin{aligned}
\|(-A)^\alpha\exp(At)\|&=\|(-A)^{\alpha-n}A^n\exp(At)\|\\
&\leq\dfrac{1}{\Gamma(n-\alpha)}\int_0^\infty s^{n-\alpha-1}\|A^n\exp(A(t+s))\|\,\mathrm{d}s\\
&\leq\dfrac{M_n}{\Gamma(n-\alpha)}\int_0^\infty s^{n-\alpha-1}(t+s)^{-n}e^{-\delta(t+s)}\,\mathrm{d}s\\
&\leq\dfrac{M_ne^{-\delta t}}{\Gamma(n-\alpha)t^\alpha}\int_0^\infty r^{n-\alpha-1}(1+r)^{-n}\,\mathrm{d}r.
\end{aligned}\]最后有\[\begin{aligned}
\|e^{At}u-u\|&=\left\|\int_0^t Ae^{As}u\,\mathrm{d}s\right\|=\left\|\int_0^t(-A)^{1-\alpha}e^{As}(-A)^\alpha u\,\mathrm{d}s\right\|\\
&\leq C\int_0^t s^{\alpha-1}\|(-A)^\alpha u\|\,\mathrm{d}s\leq C_\alpha t^\alpha\|(-A)^\alpha u\|.
\end{aligned}\]证明完成.
\end{Proof}

\subsubsection{解析半群的扰动}\label{subsec:analytic-semigroup-perturbations}

设 $A$ 为某半群的无穷小生成元. $B$ 必须在何种意义下 ``小\textquotesingle\textquotesingle, 才能保证 $A+B$ 也生成半群? 对 $C_0$ 半群, 在 Problem~\ref{ex:quasicontraction-bounded-operator-perturbation} 中已看到, $B$ 有界就足够. 实际上, 除非施加小性以外的条件, 例如 $A$ 和 $B$ 均耗散, 这基本上就是所能得到的最佳结果. 特别地, $B$ 为 ``低阶的\textquotesingle\textquotesingle\ 并无帮助. 读者可以验证, 虽然算子 $Au=u_{xxx}$ 在 $L^2(\mathbb{R})$ 中生成 $C_0$ 半群, 但算子 $(A+B)u=u_{xxx}-u_{xx}$ 并不生成. 然而, 对解析半群, 情形好得多. 主要结果为以下定理.

\begin{Theorem}
\label{thm:analytic-generator-small-relative-perturbation}
设 $A$ 为解析半群的无穷小生成元. 则存在正数 $\delta$, 使若 $B$ 为满足以下条件的任意算子:\begin{enumerate}\item $B$ 为闭算子, 且 $\mathcal{D}(B)\supseteq\mathcal{D}(A)$,\item 对 $u\in\mathcal{D}(A)$, 有 $\|Bu\|\leq a\|Au\|+b\|u\|$, 其中 $a<\delta$,\end{enumerate}则 $A+B$ 也是解析半群的无穷小生成元.
\end{Theorem}

\begin{Proof}
由于 $A$ 生成解析半群, 存在 $\omega\in\mathbb{R}$ 和 $M>0$, 使对 $\operatorname{Re}\lambda>\omega$ 有 $\|R_\lambda(A)\|\leq M/|\lambda-\omega|$. 算子 $BR_\lambda(A)$ 有界, 并且有\[\begin{aligned}
\|BR_\lambda(A)u\|&\leq a\|AR_\lambda(A)u\|+b\|R_\lambda(A)u\|\\
&\leq a\left(1+\dfrac{M|\lambda|}{|\lambda-\omega|}\right)\|u\|+\dfrac{bM}{|\lambda-\omega|}\|u\|.
\end{aligned}\]对任意 $\epsilon>0$, 可找到 $\omega^{\prime}$, 使对 $\operatorname{Re}\lambda>\omega^{\prime}$ 有 $\|BR_\lambda(A)\|\leq a(1+M+\epsilon)$. 若还有 $a<(1+M)^{-1}$, 则可选择 $\epsilon$, 使对 $\operatorname{Re}\lambda>\omega^{\prime}$ 有 $\|BR_\lambda(A)\|<1$. 其余由恒等式
\begin{equation}\label{eq:analytic-generator-perturbation-resolvent-identity}
R_\lambda(A+B)=R_\lambda(A)(I+BR_\lambda(A))^{-1},
\end{equation}
得到; 由此可知对 $\operatorname{Re}\lambda>\omega^{\prime}$ 有 $\|R_\lambda(A+B)\|\leq M^{\prime}/|\lambda-\omega^{\prime}|$.
\end{Proof}

在应用中, $B$ 通常比 $A$ ``阶数更低\textquotesingle\textquotesingle, 且前一定理中的 $a$ 可取为任意小. ``低阶\textquotesingle\textquotesingle\ 这一概念的抽象形式可用分数次幂表述. 有以下引理.

\begin{Lemma}
\label{lem:analytic-fractional-lower-order-relative-bound}
设 $A$ 为解析半群的无穷小生成元, 并假设 $B$ 为闭算子, 且对某个 $\alpha\in(0,1)$ 有 $\mathcal{D}(B)\supseteq\mathcal{D}((\omega I-A)^\alpha)$. 则存在常数 $C$, 使
\begin{equation}\label{eq:analytic-generator-fractional-lower-order-relative-bound}
\|Bu\|\leq C(\rho^\alpha\|u\|+\rho^{\alpha-1}\|(A-\omega I)u\|)
\end{equation}
对每个 $u\in\mathcal{D}(A)$ 及每个 $\rho>0$ 成立.
\end{Lemma}

选取足够大的 $\rho$ 并应用前一定理, 得到 $A+B$ 生成解析半群.

\begin{Proof}
不失一般性, 可假设 $\omega=0$. 若 $\mathcal{D}(B)\supseteq\mathcal{D}((-A)^\alpha)$, 则 $B(-A)^{-\alpha}$ 有界; 即存在常数 $C$, 使 $\|Bu\|\leq C\|(-A)^\alpha u\|$. 因而只需对 $B=(-A)^\alpha$ 证明 \eqref{eq:analytic-generator-fractional-lower-order-relative-bound}. 对 $u\in\mathcal{D}(A)$, 有\[\begin{aligned}
\|(-A)^\alpha u\|&\leq\left|\dfrac{\sin(\pi\alpha)}{\pi}\right|\int_0^\rho\lambda^{\alpha-1}\|AR_\lambda(A)\|\|u\|\,\mathrm{d}\lambda\\
&\quad+\left|\dfrac{\sin(\pi\alpha)}{\pi}\right|\int_\rho^\infty\lambda^{\alpha-1}\|R_\lambda(A)\|\|Au\|\,\mathrm{d}\lambda.
\end{aligned}\]再利用 $\|R_\lambda(A)\|\leq M/\lambda$ 和 $\|AR_\lambda(A)\|\leq1+M$, 即完成引理的证明.
\end{Proof}

在应用中, 往往很难精确刻画分数次幂的定义域. 通常并非验证 $\mathcal{D}(B)\supseteq\mathcal{D}((\omega I-A)^\alpha)$, 而是直接验证 \eqref{eq:analytic-generator-fractional-lower-order-relative-bound}. 在这种情形下, 以下结果值得关注.

\begin{Lemma}
\label{lem:analytic-relative-estimate-implies-fractional-domain}
设 $A$ 为解析半群的生成元, $B$ 为闭线性算子, 使 $\mathcal{D}(B)\supseteq\mathcal{D}(A)$, 且对某个 $\gamma\in(0,1)$ 及每个 $\rho\geq\rho_0>0$, 有
\begin{equation}\label{eq:analytic-generator-relative-estimate-fractional-domain-hypothesis}
\|Bu\|\leq C(\rho^\gamma\|u\|+\rho^{\gamma-1}\|Au\|)
\end{equation}
对每个 $u\in\mathcal{D}(A)$ 成立. 则对每个 $\alpha>\gamma$, 有 $\mathcal{D}(B)\supseteq\mathcal{D}((\omega-A)^\alpha)$.
\end{Lemma}

\begin{Proof}
再次不失一般性地假设 $\omega=0$. 设 $u\in\mathcal{D}((-A)^{1-\alpha})$, 从而 $(-A)^{-\alpha}u\in\mathcal{D}(A)$. 若积分收敛, 则有
\begin{equation}\label{eq:analytic-generator-fractional-domain-converse-integral}
B(-A)^{-\alpha}u=\dfrac{1}{\Gamma(\alpha)}\int_0^\infty t^{\alpha-1}B\exp(At)\,\mathrm{d}t,
\end{equation}
将该积分拆成
\begin{equation}\label{eq:analytic-generator-fractional-domain-converse-integral-split}
\int_0^\delta t^{\alpha-1}B\exp(At)\,\mathrm{d}t+\int_\delta^\infty t^{\alpha-1}B\exp(At)\,\mathrm{d}t.
\end{equation}
置 $\delta=1/\rho_0$, 在第二个积分中取 $\rho=\rho_0$, 在第一个积分中取 $\rho=1/t$, 并应用 \eqref{eq:analytic-generator-relative-estimate-fractional-domain-hypothesis}. 得到当 $\alpha>\gamma$ 时 $B(-A)^{-\alpha}$ 有界, 由此得到引理.
\end{Proof}

现在给出对抛物型 PDE 的一个应用. 设 $\Omega$ 为 $\mathbb{R}^m$ 中具有光滑边界的有界区域, $a_{ij}(\bm{x})$ 为 $C^1(\overline\Omega)$ 类, 且矩阵 $a_{ij}$ 对称并严格正定; 设 $b_i(\bm{x})$, $c(\bm{x})$ 为 $C^0(\overline\Omega)$ 类. 在 $L^2(\Omega)$ 中考虑算子
\begin{equation}\label{eq:analytic-semigroup-elliptic-dirichlet-generator}
Au=\dfrac{\partial}{\partial x_i}\left(a_{ij}(\bm{x})\dfrac{\partial u}{\partial x_j}\right)+b_i(\bm{x})\dfrac{\partial u}{\partial x_i}+c(\bm{x})u
\end{equation}
其定义域为 $H^2(\Omega)\cap H_0^1(\Omega)$. 断言

\begin{Theorem}
\label{thm:elliptic-dirichlet-analytic-generation}
$A$ 生成解析半群.
\end{Theorem}

\begin{Proof}
令
\begin{equation}\label{eq:analytic-semigroup-elliptic-principal-self-adjoint-part}
A_0u=\dfrac{\partial}{\partial x_i}\left(a_{ij}(\bm{x})\dfrac{\partial u}{\partial x_j}\right).
\end{equation}
则 $A_0$ 为具有负谱的自伴算子; 因而显然生成解析半群. 此外, 有\[\begin{aligned}
\|(A-A_0)u\|_2&\leq C\|u\|_{1,2}\leq C\|u\|_2^{1/2}\|u\|_{2,2}^{1/2}\\
&\leq C\|u\|_2^{1/2}\|A_0u\|_2^{1/2}\leq C(\rho^{1/2}\|u\|_2+\rho^{-1/2}\|A_0u\|_2).
\end{aligned}\]因此, 对任意 $\alpha>1/2$, $\mathcal{D}(A-A_0)$ 包含 $\mathcal{D}((-A_0)^\alpha)$.
\end{Proof}

\begin{Remark}
\label{rem:elliptic-square-root-dirichlet-domain}
聪明的读者可能猜想, $\mathcal{D}((-A_0)^{1/2})$ 实际上就是 $H_0^1(\Omega)$. 这一猜想确有根据. 然而, 其证明会比上面的讨论复杂得多.
\end{Remark}

\subsubsection{温和解的正则性}\label{subsec:analytic-mild-solution-regularity}

现在转向非齐次初值问题
\begin{equation}\label{eq:analytic-semigroup-inhomogeneous-mild-initial-value-problem}
\dot u(t)=Au(t)+f(t),\ u(0)=u_0.
\end{equation}
其温和解为
\begin{equation}\label{eq:analytic-semigroup-mild-solution-representation}
u(t)=e^{At}u_0+\int_0^t e^{A(t-s)}f(s)\,\mathrm{d}s.
\end{equation}
若 $A$ 生成解析半群, 已经知道项 $e^{At}u_0$ 对 $t>0$ 关于 $t$ 解析; 此外, $e^{At}u_0$ 对每个 $n$ 都属于 $\mathcal{D}(A^n)$. 还知道当 $t\to0$ 时, 有 $\|A^ne^{At}u_0\|\leq C\|u_0\|/t^n$. 因而只需关注项
\begin{equation}\label{eq:analytic-semigroup-mild-inhomogeneous-convolution}
v(t):=\int_0^t e^{A(t-s)}f(s)\,\mathrm{d}s.
\end{equation}
需要以下定义:

\begin{Definition}
\label{def:banach-valued-holder-continuity}
若存在常数 $L$, 使
\begin{equation}\label{eq:banach-valued-holder-continuity-definition}
\|f(t)-f(s)\|\leq L|t-s|^\theta\quad\forall s,t\in[0,T].
\end{equation}
则称 $f\in C^\theta([0,T];X)$, $0<\theta<1$.
\end{Definition}

\begin{Lemma}
\label{lem:analytic-holder-data-difference-convolution}
设 $A$ 为 $X$ 中解析半群的无穷小生成元, 并假设对某个 $\theta\in(0,1)$ 有 $f\in C^\theta([0,T];X)$. 令
\begin{equation}\label{eq:analytic-semigroup-holder-data-difference-convolution}
w(t)=\int_0^t e^{A(t-s)}(f(s)-f(t))\,\mathrm{d}s.
\end{equation}
则对每个 $t\in[0,T]$, 有 $w(t)\in\mathcal{D}(A)$, 且 $Aw\in C^\theta([0,T];X)$.
\end{Lemma}

\begin{Proof}
假设对 $t\in(0,T]$ 有 $\|\exp(At)\|\leq M$ 和 $\|A\exp(At)\|\leq C/t$. 由估计
\begin{equation}\label{eq:analytic-holder-data-difference-in-generator-domain-bound}
\left\|\int_0^t Ae^{A(t-s)}(f(s)-f(t))\,\mathrm{d}s\right\|\leq\int_0^t\dfrac{C}{t-s}L(t-s)^\theta\,\mathrm{d}s=\dfrac{CLt^\theta}{\theta}.
\end{equation}
可得 $w(t)\in\mathcal{D}(A)$. 还需证明 Hölder 连续性. 首先注意到
\begin{equation}\label{eq:analytic-generator-semigroup-time-difference-bound}
\begin{aligned}\|A\exp(At)-A\exp(As)\|&=\left\|\int_s^t A^2\exp(A\tau)\,\mathrm{d}\tau\right\|\\&\leq\int_s^t\|A^2\exp(A\tau)\|\,\mathrm{d}\tau\\&\leq4C\int_s^t\tau^{-2}\,\mathrm{d}\tau=4Ct^{-1}s^{-1}(t-s).\end{aligned}
\end{equation}
再写成\[\begin{aligned}
&Aw(t+h)-Aw(t)\\
&\quad=A\int_0^t[\exp(A(t+h-s))-\exp(A(t-s))][f(s)-f(t)]\,\mathrm{d}s\\
&\qquad+A\int_0^t\exp(A(t+h-s))(f(t)-f(t+h))\,\mathrm{d}s\\
&\qquad+A\int_t^{t+h}\exp(A(t+h-s))(f(s)-f(t+h))\,\mathrm{d}s\\
&\quad=:I_1+I_2+I_3.
\end{aligned}\]现在利用 \eqref{eq:analytic-generator-semigroup-time-difference-bound}, 得到
\begin{equation}\label{eq:analytic-holder-difference-convolution-first-term-bound}
\|I_1\|\leq\int_0^t4C(t+h-s)^{-1}(t-s)^{-1}hL(t-s)^\theta\,\mathrm{d}s\leq C_1h^\theta.
\end{equation}
(使用代换 $s=t-h\tau$.) $I_2$ 可改写为 $(\exp(A(t+h))-\exp(At))(f(t)-f(t+h))$, Hölder 估计直接由第二个因子的 Hölder 估计得到. 对最后一项, 有
\begin{equation}\label{eq:analytic-holder-difference-convolution-third-term-bound}
\|I_3\|\leq\int_t^{t+h}\dfrac{C}{t+h-s}L(t+h-s)^\theta\,\mathrm{d}s\leq C_2h^\theta.
\end{equation}

\end{Proof}

下面的正则性结果说明, 除去 $t=0$ 的某个邻域, $\dot u$ 和 $Au$ 与 $f$ 一样光滑. 注意, 对 $C_0$ 半群没有可与之相比的结果.

\begin{Theorem}
\label{thm:analytic-mild-holder-regularity}
设 $A$ 为解析半群的无穷小生成元, $u$ 为 \eqref{eq:analytic-semigroup-inhomogeneous-mild-initial-value-problem} 的由 \eqref{eq:analytic-semigroup-mild-solution-representation} 给出的解. 此外, 设 $f\in C^\theta([0,T];X)$. 则:\begin{enumerate}\item\label{concl:analytic-mild-regularity-away-from-zero} 对每个 $\delta>0$, $Au$ 和 $\dot u$ 都属于 $C^\theta([\delta,T];X)$.\item\label{concl:analytic-mild-regularity-continuity-at-zero} 若 $u_0\in\mathcal{D}(A)$, 则 $Au$ 和 $\dot u$ 都属于 $C([0,T];X)$.\item\label{concl:analytic-mild-regularity-zero-compatible-holder} 若 $u_0=0$ 且 $f(0)=0$, 则 $Au$ 和 $\dot u$ 都属于 $C^\theta([0,T];X)$.\end{enumerate}
\end{Theorem}

\begin{Proof}
设 $v(t)$ 如 \eqref{eq:analytic-semigroup-mild-inhomogeneous-convolution} 所给. 回顾 Theorem~\ref{thm:semigroup-mild-solution-classical-regularity} 的证明, 若 $Av(t)$ 存在, 则 $\dot v$ 也存在, 且 $\dot v=Av+f$.\par 因而验证该定理时只需考虑 $Au$. 为此写成
\begin{equation}\label{eq:analytic-mild-regularity-three-part-solution-splitting}
u(t)=e^{At}u_0+\int_0^t e^{A(t-s)}(f(s)-f(t))\,\mathrm{d}s+\int_0^t e^{A(t-s)}f(t)\,\mathrm{d}s.
\end{equation}
由前一引理, 只需考虑最后一项. 有
\begin{equation}\label{eq:analytic-mild-regularity-frozen-data-generator-integral}
A\int_0^t e^{A(t-s)}f(t)\,\mathrm{d}s=(e^{At}-I)f(t).
\end{equation}
由于已假设 $f$ 属于 $C^\theta([0,T];X)$, 只需考虑 $\exp(At)f(t)$. 对 $t\geq\delta$ 和 $h>0$, 有
\begin{equation}\label{eq:analytic-mild-regularity-holder-estimate-away-from-zero}
\begin{aligned}[t]&\|\exp(A(t+h))f(t+h)-\exp(At)f(t)\|\\&\quad\leq\|\exp(A(t+h))\|\ \|f(t+h)-f(t)\|\\&\qquad+\|\exp(A(t+h))-\exp(At)\|\ \|f(t)\|\\&\quad\leq C_1h^\theta+C_2\dfrac{h}{\delta}.\end{aligned}
\end{equation}
这就给出第~\ref{concl:analytic-mild-regularity-away-from-zero} 部分. 再注意到
\begin{equation}\label{eq:analytic-mild-regularity-initial-strong-continuity-bound}
\begin{aligned}\|\exp(At)f(t)-f(0)\|&\leq\|\exp(At)f(0)-f(0)\|\\&\quad+\|\exp(At)\|\|f(t)-f(0)\|;\end{aligned}
\end{equation}
由半群的强连续性即得第~\ref{concl:analytic-mild-regularity-continuity-at-zero} 部分. 为证明第~\ref{concl:analytic-mild-regularity-zero-compatible-holder} 部分, 先按 \eqref{eq:analytic-mild-regularity-holder-estimate-away-from-zero} 的方式处理, 然后估计\[\begin{aligned}
\|(e^{A(t+h)}-e^{At})f(t)\|&=\left\|\int_t^{t+h}Ae^{A\tau}f(t)\,\mathrm{d}\tau\right\|\\
&\leq\int_t^{t+h}\|Ae^{A\tau}(f(t)-f(0))\|\,\mathrm{d}\tau\\
&\leq C\int_t^{t+h}\tau^{-1}t^\theta\,\mathrm{d}\tau\\
&\leq C\int_t^{t+h}\tau^{\theta-1}\,\mathrm{d}\tau\leq Ch^\theta.
\end{aligned}\]证明完成
\end{Proof}

\Needspace{5\baselineskip}\textbf{Problems}

\begin{enumerate}[label=\thesection.\arabic*.,ref=\thesection.\arabic*,start=20,series=book-problems]

\item\label{ex:analytic-semigroup-derivative-bound-proof} 证明 Lemma~\ref{lem:analytic-semigroup-time-derivative-bound}.

\item\label{ex:analytic-semigroup-spectral-growth-proof} 证明 Theorem~\ref{thm:analytic-semigroup-spectrum-controls-growth}.

\item\label{ex:fractional-power-integer-inverse-identity} 验证对 $n\in\mathbb{N}$, 有 $(-A)^{-n}=(-A^{-1})^n$, 其中 $(-A)^{-n}$ 由 \eqref{eq:analytic-generator-negative-fractional-power-contour} 定义.

\item\label{ex:analytic-mild-fractional-compatible-initial-regularity} 假设 $u_0$ 和 $f(0)$ 属于 $\omega-A$ 的适当分数次幂的定义域, 证明 Theorem~\ref{thm:analytic-mild-holder-regularity} 的第~\ref{concl:analytic-mild-regularity-zero-compatible-holder} 部分.

\item\label{ex:analytic-semigroup-off-diagonal-product-space-generator} 设 $A$ 为 $X$ 上解析半群的无穷小生成元, $B$ 为满足 $\mathcal{D}(B)\supseteq\mathcal{D}(A)$ 的闭算子. 证明由 $\mathcal{A}(u,v)=(Bv,Av)$ 定义的算子在 $X\times X$ 上生成解析半群.

\item\label{ex:analytic-semigroup-strongly-damped-wave-equation} 讨论如何将解析半群应用于带 Dirichlet 边界条件的方程 $u_{tt}=\Delta u_t+\Delta u$.

\end{enumerate}
